\pdfminorversion=6
\documentclass[aspectratio=169]{beamer}
\usepackage{geometry}
\geometry{margin=0.8cm}
\usepackage{graphicx}

\mode<presentation>
{
  % Tema elegante con fuentes Serif para un estilo académico y clásico
  \usetheme{Boadilla}
  \usecolortheme{rose} 
  \usefonttheme{serif} 
  
  \setbeamertemplate{caption}[numbered]
  \setbeamertemplate{navigation symbols}{} % Limpia la pantalla de iconos de navegación
  \setbeamertemplate{footline}[frame number] 
} 

% Definición de comandos adaptados al nuevo tema
\newcommand{\saga}{\textbf{PSICOLOGÍA DEL COLOR}}
\newcommand{\yo}{Diana Anahi Buenavides Martinez}
\newcommand{\magia}[1]{\textit{\textbf{\color{burgundy}#1}}}
\newcommand{\escuela}{Facultad de Estudios Profesionales Zona Media}

\usepackage[spanish]{babel}
\usepackage[utf8]{inputenc} % Cambiado a utf8 estándar para evitar conflictos en Overleaf
\usepackage{tikz}
\usepackage{courier}
\usepackage{array}
\usepackage{bold-extra}
\usepackage[thicklines]{cancel}
\usepackage{fancyvrb}
\usepackage{colortbl} % Necesario para colorear filas de tablas sin errores

% Paleta de colores basada en la Psicología del Color de Hogwarts
\xdefinecolor{burgundy}{rgb}{0.55, 0.05, 0.15}      % Rojo Gryffindor: Pasión, valor, impulso
\xdefinecolor{wizardblue}{rgb}{0.12, 0.29, 0.59}   % Azul Ravenclaw: Inteligencia, calma, lógica
\xdefinecolor{serpentgreen}{rgb}{0.10, 0.40, 0.25} % Verde Slytherin: Ambición, poder, astucia
\xdefinecolor{oldgold}{rgb}{0.78, 0.60, 0.15}      % Amarillo/Oro Hufflepuff: Lealtad, calidez, esfuerzo
\xdefinecolor{charcoal}{rgb}{0.20, 0.20, 0.22}     % Gris oscuro para neutralidad y estructura
\xdefinecolor{softcream}{rgb}{0.98, 0.97, 0.94}    % Fondo pergamino suave para bloques

% Ajustamos los colores globales de Beamer
\setbeamercolor{palette primary}{bg=charcoal, fg=white}
\setbeamercolor{palette secondary}{bg=oldgold, fg=white}
\setbeamercolor{palette tertiary}{bg=burgundy, fg=white}
\setbeamercolor{titlelike}{bg=charcoal, fg=white}
\setbeamercolor{structure}{fg=burgundy} 
\setbeamercolor{block title}{bg=charcoal!90, fg=white}
\setbeamercolor{block body}{bg=softcream, fg=black}

\title{\saga}
\subtitle{El Significado de la Identidad Visual en Hogwarts}
\author{\yo}
\institute{\escuela}
\date{\today}

\begin{document}

% Barra superior elegante con el lema clásico
\logo{\pgfputat{\pgfxy(0.11, 7.4)}{\pgfbox[right,base]{\tikz{
    \filldraw[fill=charcoal, draw=none] (0 cm, 0.1 cm) rectangle (50 cm, 1 cm);
    \filldraw[fill=oldgold, draw=none] (0 cm, 0 cm) rectangle (50 cm, 0.1 cm);
}\mbox{\hspace{-9 cm}\small\color{oldgold}\fontfamily{pzc}\selectfont ~~~~~~Draco Dormiens Nunquam Titillandus}}}}

% Diapositiva 1: Portada
\begin{frame}
  \titlepage
\end{frame}

% Diapositiva 2: Introducción
\begin{frame}{\saga}
\vspace{0.1 cm}

\begin{block}{Introducción: La Fuerza de la Identidad Visual}
{\fontsize{9.5pt}{11pt}\selectfont 
El color no es meramente estético; genera un impacto psicológico inmediato en el espectador. J.K. Rowling y los diseñadores de la saga utilizaron la psicología del color para segmentar las personalidades, virtudes y defectos de las cuatro facciones del Colegio Hogwarts, convirtiendo los uniformes en un reflejo del alma de sus integrantes.
}
\end{block}

\vspace{-0.1 cm}
{\fontsize{9.5pt}{11pt}\selectfont
\begin{itemize}
  \item \textbf{Cromatismo Emocional:} Los colores activan respuestas inconscientes asociadas al peligro, la calma, la ambición o la pureza.
  \item \textbf{La Dualidad del Color:} Cada tono utilizado posee una cara positiva (virtud) y una negativa (defecto o extremismo).
  \item \textbf{Narrativa Visual:} El espectador identifica las intenciones y el trasfondo de un personaje con tan solo ver su heráldica.
\end{itemize}
}
\end{frame}

% Diapositiva 3: El Trío de Oro y su espectro
\begin{frame}{El Impacto Teórico}
\begin{center}
    \large\color{burgundy}\textbf{La Tríada del Comportamiento Humano}
    \vfill
    \begin{minipage}{0.75\linewidth}
    \centering\itshape
    "La selección de colores en la literatura de fantasía actúa como un ancla psicológica. El rojo despierta la acción inmediata, el azul la reflexión sesuda y el amarillo la constancia incondicional."
    \end{minipage}
    \vfill
    \tikz{
        \draw[line width=1.5pt, color=oldgold] (-2,0) -- (2,0);
        \node[draw=burgundy, fill=softcream, thick, shape=circle, inner sep=4pt] at (0,0) {\color{burgundy}\textbf{P.C}};
    }
\end{center}
\end{frame}

% Diapositiva 4: El Castillo y el Lienzo Neutral
\begin{frame}{La Armonía del Entorno}
\vfill
\begin{center}
    \Large \color{charcoal}\textbf{\textit{"El color comunica allí donde las palabras no llegan."}}
    \vfill
    \begin{block}{El Contraste de las Estancias}
    Mientras las áreas comunes del castillo mantienen tonos grises y pétreos (neutralidad y sabiduría antigua), las salas comunes rompen la monotonía utilizando la saturación cromática para reafirmar la identidad de sus estudiantes.
    \end{block}
\end{center}
\vfill
\end{frame}

% Diapositiva 5: Gráfico de las Cuatro Casas (Uso de TikZ mejorado)
\begin{frame}{El Mapa Psicológico de las Casas}
    \begin{center}
    \small\color{charcoal} El equilibrio del colegio radica en la compensación de estas cuatro fuerzas cromáticas:
    \vfill
    \begin{tikzpicture}[node distance=3.2cm, every node/.style={font=\sffamily\bfseries\footnotesize, inner sep=8pt}]
        \node (gry) [draw=burgundy, fill=burgundy!15, text=burgundy, rounded corners=3pt] {Gryffindor (Pasión/Valor)};
        \node (rav) [draw=wizardblue, fill=wizardblue!15, text=wizardblue, rounded corners=3pt, right of=gry] {Ravenclaw (Mente/Lógica)};
        \node (sly) [draw=serpentgreen, fill=serpentgreen!15, text=serpentgreen, rounded corners=3pt, below of=rav, node distance=2.2cm] {Slytherin (Poder/Astucia)};
        \node (huf) [draw=oldgold, fill=oldgold!20, text=black, rounded corners=3pt, left of=sly] {Hufflepuff (Lealtad/Paz)};
        
        \draw[<->, line width=1.2pt, color=charcoal!40] (gry) -- (rav);
        \draw[<->, line width=1.2pt, color=charcoal!40] (rav) -- (sly);
        \draw[<->, line width=1.2pt, color=charcoal!40] (sly) -- (huf);
        \draw[<->, line width=1.2pt, color=charcoal!40] (huf) -- (gry);
    \end{tikzpicture}
    \end{center}
\end{frame}

% Diapositiva 6: Tabla Informativa
\begin{frame}{Significados Psicológicos según la Casa}
   \begin{center}
   \small Análisis del lenguaje del color aplicado a las virtudes principales:
   \vfill
   \begin{tabular}{|l|c|l|}
   \hline
   \rowcolor{charcoal} \color{white}\textbf{Casa} & \color{white}\textbf{Tono Predominante} & \color{white}\textbf{Efecto Psicológico Asociado} \\ \hline
   \textbf{Gryffindor} & Rojo Granate & Estimula la adrenalina, el coraje, la osadía y el peligro. \\ \hline
   \textbf{Ravenclaw}  & Azul Profundo & Fomenta la concentración, la introspección, la paz y la lógica. \\ \hline
   \textbf{Slytherin}  & Verde Esmeralda & Evoca el estatus, la ambición, el misterio y el renacimiento. \\ \hline
   \textbf{Hufflepuff} & Amarillo Oro & Transmite optimismo, calidez, trabajo en equipo y constancia. \\ \hline
   \end{tabular}
   \end{center}
\end{frame}

% Diapositiva 7: Disciplinas y su color
\begin{frame}{La Percepción de la Magia}
    \begin{block}{Efectos Visuales en el Espectador}
    La psicología del color también segmenta los tipos de magia mostrados a lo largo de la historia:
    \end{block}
    \vfill
    \begin{columns}[T]
        \begin{column}{0.48\linewidth}
            \begin{itemize}
                \item \magia{Chispazos Rojos / Dorados:} Asociados a la protección, la defensa activa y los hechizos de desarme benévolos.
                \item \magia{Auras Azules / Blancas:} Destinados a la alta magia cognitiva, encantamientos de memoria y proyecciones mentales.
            \end{itemize}
        \end{column}
        \begin{column}{0.48\linewidth}
            \begin{itemize}
                \item \magia{Destellos Verdes Oscuros:} Relacionados intrínsecamente con la corrupción, la ambición desmedida y la letalidad.
                \item \magia{Luminosidad Cálida:} Utilizada en entornos seguros, de sanación, herbolaria y confort hogareño.
            \end{itemize}
        \end{column}
    \end{columns}
\end{frame}

% Diapositiva 8: Conclusión
\begin{frame}{Conclusión}
\begin{center}
    \tikz\draw[line width=1pt, color=oldgold] (-3,0) -- (3,0);
    \vfill
    \begin{minipage}{0.85\linewidth}
    \centering\fontsize{11pt}{14pt}\selectfont
    El uso del color en la narrativa demuestra que la identidad visual es una herramienta indispensable para conectar con el inconsciente del público. 
    
    Al equilibrar las cuatro casas mediante paletas cromáticas opuestas pero complementarias, se nos enseña que ninguna virtud es autosuficiente: la mente necesita del valor, y la ambición requiere de la lealtad para no perder su rumbo.
    \end{minipage}
    \vfill
    \tikz\draw[line width=1pt, color=oldgold] (-3,0) -- (3,0);
\end{center}
\end{frame}

\end{document}